% Lösungen Einheit 3 · Summe mal Summe
\zweigkopf{Einheit 3 von 5 · Summe mal Summe}{Hier lernst du, Terme mit zwei Variablen zusammenzufassen und zwei Klammern miteinander malzunehmen · neu oder schon bekannt – je nach Buch · keine P10-Aufgabe · baut auf: Klammern auflösen (Einheit 1), Terme zusammenfassen, Terme malnehmen}{l3}
\erg{20}{je vier Bögen: a) $x \cdot x$, $x \cdot 9$, $1 \cdot x$, $1 \cdot 9$ \quad b) $m \cdot n$, $m \cdot 4$, $3 \cdot n$, $3 \cdot 4$ \quad c) $x \cdot x$, $x \cdot 1$, $(-5) \cdot x$, $(-5) \cdot 1$ \quad d) $2y \cdot y$, $2y \cdot (-7)$, $1 \cdot y$, $1 \cdot (-7)$}
\erg{21}{a) $5$ \quad b) $7$ \quad c) $4$ \quad d) $17$ \quad e) $1$ \quad f) $11$ \quad g) $-12$ \quad h) $19$}
\erg{22}{a) $7x + 2y$ \quad b) $5m + 6n$ \quad c) $8m + 7n$ \quad d) $7u + 7v$ \quad e) $7x + 3y$ \quad f) $3x^2 + 2xy + x$ \quad g) $3mn + 4m$ \quad h) $6x + 3y + 4x = 10x + 3y$}
\erg{23}{a) $x^2 + 4x + x + 4$ \quad b) $m^2 + 7m + 2m + 14$ \quad c) $y^2 + y + 3y + 3$ \quad d) $uv + 3u + 4v + 12$ \quad e) $x^2 + 9x + 2x + 18 = x^2 + 11x + 18$}
\erg{24}{a) $x^2 - 4x - 12$ \quad b) $x^2 - 10x + 21$ \quad c) $2x^2 + 11x + 5$ \quad d) $3m^2 + 7mn + 2n^2$ \quad e) $8x^2 - 20x - 6x + 15 = 8x^2 - 26x + 15$}
\erg{25}{a) Paul hat nur Erstes mal Erstes und Letztes mal Letztes gerechnet, $6x$ und $4x$ fehlen; richtig: $x^2 + 4x + 6x + 24 = x^2 + 10x + 24$ \quad b) $x^2 + 9x + 14$ \quad c) Vorzeichen: $(-5) \cdot (-1) = +5$; richtig: $x^2 - 6x + 5$ \quad d) $x^2 - 10x + 16$ \quad e) Lukas hat $x^2$ mit $x$-Gliedern zusammengefasst; richtig: $x^2 + 6x + 8$ \quad f) $y^2 + 9y + 18$}
\erg{26}{a) $x^2$, $4x$ (oben rechts), $x$ (unten links), $4$ (unten rechts) \quad b) Die Seiten bestehen aus je zwei Stücken; jedes Stück der einen Seite bildet mit jedem Stück der anderen ein Teilrechteck – das sind $2 \cdot 2 = 4$ Teilrechtecke, also vier Produkte. \quad c) Er hat nur $x^2$ und $4$; es fehlen die beiden Teilrechtecke $4x$ und $x$. Richtig: $x^2 + 5x + 4$}
\erg{27}{a) $(s + 4) \cdot (s - 2) = s^2 - 2s + 4s - 8 = s^2 + 2s - 8$ \quad b) alt $36$~m², neu $6^2 + 2 \cdot 6 - 8 = 40$~m² \quad c) Unterschied: $s^2 + 2s - 8 - s^2 = 2s - 8$; größer, wenn $2s - 8 > 0$, also für $s > 4$ (bei $s = 4$ gleich groß, für $2 < s < 4$ kleiner)}
